\documentclass[sigconf]{acmart}

\usepackage{amsmath}

\AtBeginDocument{%
  \providecommand\BibTeX{{Bib\TeX}}}

\setcopyright{none}
\settopmatter{printacmref=false}

\begin{document}

\title{Does Model Capacity Improve Candidate Ranking Quality? 
An Evaluation of Pretrained Audio Representations for Music Recommendation in Classical Archives}

\author{Keita Katsumi}
\affiliation{%
  \institution{California State Polytechnic University, Pomona}
  \city{Pomona}
  \state{California}
  \country{USA}
}
\email{kkatsumi@cpp.edu}

\renewcommand{\shortauthors}{Katsumi}

\begin{abstract}
In cold-start recommendation settings where user interaction data are unavailable, candidate generation relies heavily on content-based representations. This study investigates the effect of embedding-model capacity on candidate ranking quality in such a setting, using a small-scale classical music archive. We evaluate CNN-based and Transformer-based embedding models across capacity tiers using composer-based and musical-character retrieval tasks, measured with standard ranking and set-based IR metrics. The results show that increasing model capacity does not consistently improve ranking performance. Scaling behavior is non-monotonic across architectures, and its effect is strongly task-dependent: structured retrieval tasks exhibit greater sensitivity to capacity differences, while abstract semantic retrieval shows minimal variation across model sizes. In addition, when the number of relevant items per query is large relative to the evaluation cutoff, set-based metrics such as Recall@K and F1@K remain low despite reasonable ranking quality, whereas ranking-based metrics such as NDCG@K and Hit@K provide more stable signals. These findings suggest that, in cold-start recommendation settings without user interaction data, increasing model capacity incurs substantially higher embedding extraction cost without delivering consistent ranking improvements, and that effective candidate generation should favor lightweight or mid-sized models that achieve comparable ranking quality at substantially lower cost.
\end{abstract}

\maketitle

% Section 1: Introduction
\section{Introduction}\label{sec:introduction}
Modern music streaming platforms rely heavily on machine learning to support search, recommendation, and discovery. Much of the prior work in this area, as well as many industrial systems, is built around large-scale collaborative filtering and deep models trained on extensive user interaction logs. While effective for mainstream catalogs, these approaches are ill-suited to small, curated archives that lack dense user data and whose mission emphasizes exploration over pure popularity. Classical music archives such as the iPalpiti collection exemplify this setting: they contain high-quality, stylistically coherent recordings but limited usage data, and they require tools that can surface musically meaningful relationships among works, performances, and artists.

In such environments, content-based retrieval driven by audio embeddings becomes the primary mechanism for candidate generation. Recent advances in pretrained audio models---including lightweight convolutional networks and large transformer-based architectures---have expanded the range of available representations, but they vary substantially in parameter count and computational cost. On large-scale, heterogeneous benchmarks, increased model capacity often correlates with improved performance. Whether this trend justifies the additional cost in small, single-domain archives is less clear. In constrained archival settings, where catalog size is limited and user interaction data are absent, deploying a larger model may not improve retrieval quality enough to offset the increase in embedding extraction overhead. This raises a system design question: \textit{is scaling embedding-model capacity a cost-effective choice for content-based candidate generation in resource-constrained archival systems?}

At the same time, evaluating such systems is inherently challenging, as relevance is often broadly distributed across many items rather than concentrated in a small set of targets. This raises an additional question: do commonly used evaluation metrics reliably reflect ranking quality in these settings?

This paper addresses these questions by isolating and evaluating the backend component responsible for mapping audio recordings to embeddings and supporting similarity-based retrieval, independent of user modeling. We evaluate CNN-based and Transformer-based pretrained audio models across discrete capacity tiers, measuring both ranking quality and embedding extraction cost. Rather than assuming that larger models yield better retrieval performance, we treat model selection as a cost--quality trade-off decision and examine whether increased capacity produces gains sufficient to justify its overhead.

Two proxy retrieval tasks operationalize performance evaluation: a structured task based on composer identity and an abstract semantic task based on musical character (e.g., \emph{energetic} vs.\ \emph{calm}). These tasks probe different aspects of the embedding space and allow us to examine whether capacity benefits are uniform or task-dependent.

Ranking quality is measured using both ranking-based metrics (NDCG@K, Precision@K) and set-based metrics (Recall@K, F1@K). We further examine how these metric families behave when relevant items are broadly distributed relative to the evaluation cutoff. This research is guided by two research questions:

\noindent\textbf{RQ1:} How does embedding-model capacity influence ranking quality in content-based retrieval for a small-scale classical archive?

\noindent\textbf{RQ2:} Given the associated computational cost, under what conditions do higher-capacity models provide sufficient ranking gains to justify their deployment over lighter alternatives?

\noindent\textbf{Hypothesis:} In a small-scale, single-domain classical music archive, increasing embedding-model capacity does not consistently improve candidate ranking quality, and the performance differences that do emerge are insufficient to justify the substantially higher extraction cost.

This work makes three main contributions. First, it frames embedding-model selection as a cost--quality trade-off and provides a systematic evaluation of pretrained audio representations across capacity tiers in a realistic cold-start archival retrieval setting. Second, it demonstrates that capacity scaling is non-monotonic and task-dependent: structured retrieval based on composer identity exhibits greater sensitivity to capacity differences, while abstract semantic retrieval shows minimal variation across model sizes, and in neither case does the largest model deliver consistent gains. Third, it shows that set-based metrics such as Recall@K and F1@K can be systematically misleading when relevant items are broadly distributed, and that ranking-based metrics provide more reliable signals for evaluating content-based candidate generation quality.


% Section 2: Related Work
\section{Related Work}\label{sec:related-work}

When user interaction data are sparse or unavailable, collaborative filtering signals are weak and cold-start effects become severe \cite{schedl2018}. Content-based representations derived from the audio signal itself provide a practical basis for similarity search and candidate generation in such settings \cite{deldjoo2024}. Recent work has demonstrated the effectiveness of pretrained audio representations for music retrieval while also identifying a gap in their systematic evaluation within recommender system pipelines \cite{tamm2024}. In retrieval-based recommendation, embedding models function as candidate generation components: they map items to a representation space and rank candidates by similarity before any user-aware re-ranking is applied. The intrinsic quality of these embeddings is therefore a critical system design factor, yet the impact of embedding model capacity on candidate generation quality under data-constrained conditions remains poorly understood.

Audio embedding architectures span a wide design space, including CNNs, RNNs, Transformers, and hybrid variants \cite{zaman2023}. CNN-based architectures have demonstrated strong performance across music detection and audio classification tasks and scale effectively with increased model depth \cite{kong2020panns, jiang2024}. Transformer-based models such as AST apply self-attention directly to spectrograms and model longer-range structure \cite{gong2021ast, zaman2023}. Hybrid approaches, including CNN--RNN models for temporal modeling \cite{lin2024} and CNN--Transformer models for combining local and global representations \cite{pourmoazemi2024}, have also been explored but are treated here as background. This study focuses on pure CNN-based and pure Transformer-based families as representative architectural design choices.

While larger models often achieve stronger performance on large-scale, diverse benchmarks, their advantages in small, domain-constrained archives remain unclear. Higher-capacity models introduce substantially greater embedding extraction cost, and it is not evident that this cost is offset by meaningful gains in candidate ranking quality under cold-start, data-constrained conditions \cite{schedl2018, deldjoo2024, tamm2024}. This cost--quality trade-off is the central system design question addressed by this work.

Proper offline evaluation is a recognized challenge in music recommender systems \cite{schedl2018}. Classification benchmarks are commonly used to assess audio embeddings \cite{gong2021ast, zaman2023}, but strong classification accuracy does not necessarily imply well-structured neighborhoods for similarity-based retrieval \cite{tamm2024}. Ranking metrics such as NDCG and Precision@K are more appropriate for retrieval-oriented settings because they directly evaluate the ordering of relevant items \cite{krichene2020}. Evaluation metrics are also sensitive to relevance distribution, cutoff choices, and candidate set design; in particular, set-based metrics such as Recall@K can underestimate ranking quality when relevant items are broadly distributed relative to the cutoff \cite{krichene2020, canamares2020, urbano2013}. These issues are especially acute in small archival collections. Diversity and long-tail considerations remain relevant background concerns in settings where exploration is valued \cite{porcaro2022}.

Broader recommender-system work addresses personalization through user modeling, sequential interaction modeling, and contextual signals \cite{lin2025, abbattista2024, schedl2021}, but those concerns are outside the scope of this study. This work isolates the candidate generation stage, evaluating embedding quality independently of downstream ranking and user modeling components.



% Section 3: Method
\section{Method}\label{sec:method}\label{ch:problem}\label{ch:methodology}
This section formalizes the retrieval setting and experimental framework used to evaluate the effect of embedding-model capacity on similarity-based retrieval performance.

\subsection{Problem Formulation}

Let $D = \{x_1, x_2, \dots, x_N\}$ denote a collection of $N$ classical music recordings drawn from
the iPalpiti archive. Each item $x_i$ represents a complete audio track.

An embedding model is defined as a function:
\[
f_\theta : \mathcal{X} \rightarrow \mathbb{R}^d
\]
where $\theta$ denotes the model parameters and $d$ the embedding dimensionality. The embedding of
item $x_i$ is:
\[
z_i = f_\theta(x_i)
\]

Similarity between two items is computed using a similarity function $s(\cdot,\cdot)$, such as cosine
similarity:
\[
s(z_q, z_i) = \frac{z_q \cdot z_i}{\|z_q\|\|z_i\|}
\]
Higher similarity values indicate greater embedding-space proximity.

Given a query item $x_q \in D$, all items in the collection except $x_q$ are ranked in descending
order of similarity $s(z_q, z_i)$. Retrieval quality is evaluated with respect to a relevance
function $r_q(x_i) \in \{0,1\}$ defined by proxy tasks, and ranking quality is measured using
standard IR metrics including NDCG@K, Precision@K, Recall@K, and F1@K.

The primary factor of interest is embedding-model capacity within each architectural family. Capacity is operationalized through small, medium, and large model variants; architectural family (CNN vs. Transformer) is treated as a separate descriptive factor. Capacity effects are examined within families under a controlled retrieval setting; cross-family comparisons are interpreted descriptively.

Relevance is defined through two metadata-based proxy tasks. Let $g(x)$ denote a metadata labeling
function.

\textbf{Sanity Proxy.} Relevance is defined by shared composer identity:
\[
r_q(x_i) = 1 \quad \text{if} \quad g_{\text{composer}}(x_q) = g_{\text{composer}}(x_i)
\]

\textbf{Primary Musical Proxy.} Relevance is defined by shared character labels:
\[
r_q(x_i) = 1 \quad \text{if} \quad g_{\text{character}}(x_q) \cap g_{\text{character}}(x_i) \neq \emptyset
\]
where $g_{\text{character}}(x)$ denotes the set of character labels associated with item $x$.
This formulation allows multiple items to be relevant for a given query, reflecting the multi-label nature of semantic similarity in music.

For experimental tractability, relevance is treated as binary by projecting multi-label annotations into a match-based criterion:
\[
r_q(x_i) \in \{0,1\}
\]
where $r_q(x_i)=1$ if at least one label is shared between $x_q$ and $x_i$, and $r_q(x_i)=0$ otherwise. This formulation simplifies evaluation while preserving the underlying multi-label structure of the data, allowing multiple items to be considered relevant for a given query and resulting in a broader distribution of relevant items in the candidate set. 
This study assumes that such metadata-derived labels provide a controlled, though coarse, approximation of musical similarity.

\subsection{Retrieval Pipeline}
The backend evaluation pipeline consists of four stages: audio preprocessing, embedding extraction, similarity computation, and ranking. All components other than embedding extraction are held constant to isolate the effect of embedding capacity. The dataset is the iPalpiti classical music archive, comprising 203 tracks and 24.9 hours of audio. Ground truth is derived from editorial metadata, including composer identity and character labels used in the proxy tasks.

Each complete track serves as the retrieval unit. Because classical recordings vary substantially in duration, each track is segmented into non-overlapping 30-second chunks during offline feature extraction. Segment-level embeddings are extracted independently and then mean-pooled into a single track-level representation. This segmentation-and-pooling strategy reflects a realistic catalog-ingestion pipeline for long-form audio items and ensures that embedding extraction is feasible across all model architectures regardless of their input constraints. Cosine similarity over the resulting track-level embeddings is used for ranking across all models. No downstream personalization, collaborative signal, or learned re-ranking component is applied.

\subsection{Model Setup}
Models are selected to represent discrete capacity tiers within two architectural families: CNN-based and Transformer-based models. Within each family, small, medium, and large variants enable controlled scaling comparisons, while cross-family comparisons are interpreted descriptively because architecture, parameter count, and pretraining objectives may co-vary.

CNN-based models are drawn from a unified convolutional lineage in which scaling is achieved primarily through increased depth under comparable training objectives. Specifically, we use Cnn6 (4.8\,M parameters; small), Cnn10 (5.2\,M; medium), and Cnn14 (80.8\,M; large), following the PANNs framework proposed by Kong et al. \cite{kong2020panns}, to represent controlled capacity scaling within the CNN family.

Transformer-based models are selected from a self-supervised audio representation framework pretrained on large-scale music data. We use MERT-95M (95\,M parameters; medium) and MERT-330M (330\,M; large) to represent different capacity levels within the transformer family, following the MERT framework \cite{li2023mert}. We note that, unlike the CNN family, a small-capacity variant is not included for the Transformer models. This is because publicly available pretrained checkpoints in the MERT framework are provided solely at medium and large scales; as a result, our Transformer-based evaluation is restricted to these available variants.

Parameter counts are not strictly aligned across architectural families; therefore, cross-family comparisons are interpreted descriptively rather than as direct capacity-matched evaluations. Within each family, however, the stated parameter counts provide an explicit quantification of the capacity dimension under study. All selected models provide pretrained checkpoints and support standardized embedding extraction.

\subsection{Label Construction}

Relevance labels for the proxy tasks are derived from metadata and a semi-automatic annotation pipeline. For the sanity proxy, composer labels are normalized to a canonical form to ensure consistent matching across recordings. This normalization is performed during pseudo-label generation and stored as a stable evaluation field. For the primary musical proxy, character labels are constructed using a semi-automatic pipeline. Initial label candidates are defined based on the arousal--valence framework \cite{eerola2011}, which represents emotion in a continuous dimensional space. To support this process, pretrained music emotion recognition models are used only as auxiliary tools for generating initial label candidates. These models reflect recent advances in deep learning-based music emotion recognition \cite{jiang2024}, but are not used as a source of ground-truth relevance. Instead, the generated labels are treated as fixed proxy annotations, independent of the embedding models evaluated in this study and not aligned with their training objectives. Ambiguous cases are manually reviewed by the authors, who have domain familiarity with classical music, with disagreements resolved through discussion to ensure consistency with perceived musical characteristics. 

The resulting pseudo-labels are fixed prior to evaluation and reused across all experiments. They do not constitute ground-truth semantic similarity, but provide a controlled and reproducible basis for relative comparison across models. Because the labeling process is independent of the embedding models under study, this design isolates the effect of embedding representations without introducing model-dependent labeling bias.

\subsection{Justification of Proxy Tasks}
In practical recommender systems, evaluation is often based on historical user interaction logs. However, such signals are unavailable in the iPalpiti archive.
We therefore adopt metadata-based proxy tasks based on composer identity and musical character as deterministic offline relevance criteria. Although these proxies do not fully capture real-world user preference, they provide a controlled and reproducible basis for comparing how well different embedding models support music-similarity retrieval.

\subsection{Evaluation Protocol}

Each track in the archive is used as a query in a leave-one-out retrieval setting, where the query item itself is excluded from the candidate pool. Relevance is determined by metadata matching as defined above, and performance metrics are averaged across all query instances. Because relevance labels are binary and the dataset is relatively small, ranking ties may occur when multiple relevant items share identical similarity scores. To ensure deterministic evaluation, a consistent tie-breaking strategy (e.g., lexicographic ordering by track identifier) is applied across all models. Performance is primarily assessed using NDCG@5 and Precision@5, with interpretation focused on consistent scaling patterns rather than small numerical differences.

Furthermore, small numerical differences in mean scores do not necessarily correspond to meaningful differences in retrieval quality \cite{urbano2013}. This evaluation setup follows standard offline ranking evaluation practices in recommender systems when user interaction logs are unavailable.



% Section 4: Results
\section{Results}\label{sec:results}\label{ch:results}
\subsection{Sanity Proxy Results}

The Sanity Proxy task evaluates the ability of embedding models to retrieve tracks sharing the same composer. This task serves as a baseline for verifying that models capture fundamental compositional structures.

Table~\ref{tab:sanity_results} presents retrieval performance grouped by architectural family and ordered by capacity tier within each family.

\begin{table}[!htbp]
\centering
\caption{Sanity Proxy Results (Composer Retrieval). Metrics reported are corpus-level mean values at rank 5.}
\label{tab:sanity_results}
\resizebox{\columnwidth}{!}{%
\begin{tabular}{lccccc}
\hline
\textbf{Model} & \textbf{NDCG@5} & \textbf{Hit@5} & \textbf{Precision@5} & \textbf{Recall@5} & \textbf{F1@5} \\
\hline
\multicolumn{6}{c}{\textbf{CNN Family}} \\
\hline
CNN-Small & 0.548 & 0.640 & 0.266 & 0.195 & 0.188 \\
CNN-Medium & 0.545 & 0.640 & 0.303 & 0.215 & 0.214 \\
CNN-Large & 0.585 & 0.660 & 0.321 & 0.228 & 0.229 \\
\hline
\multicolumn{6}{c}{\textbf{Transformer Family}} \\
\hline
Transformer-Medium & 0.642 & 0.709 & 0.361 & 0.263 & 0.265 \\
Transformer-Large & 0.588 & 0.665 & 0.304 & 0.233 & 0.224 \\
\hline
\end{tabular}
}
\end{table}

The Sanity Proxy results show that capacity scaling does not produce a consistent improvement in ranking quality across architectures. Within the CNN family, NDCG@5 decreases slightly from the small to the medium model before increasing at the large tier, indicating a non-monotonic scaling pattern. Within the Transformer family, performance decreases from the medium to the large model across all metrics. These results indicate that the effect of capacity scaling is architecture-dependent and does not follow a uniform increasing trend.

Table~\ref{tab:musical_results} summarizes the performance across all models.

\begin{table}[!htbp]
\centering
\caption{Primary Musical Proxy Results (Character Retrieval). Metrics reported are corpus-level mean values at rank 5.}
\label{tab:musical_results}
\resizebox{\columnwidth}{!}{%
\begin{tabular}{lccccc}
\hline
\textbf{Model} & \textbf{NDCG@5} & \textbf{Hit@5} & \textbf{Precision@5} & \textbf{Recall@5} & \textbf{F1@5} \\
\hline
\multicolumn{6}{c}{\textbf{CNN Family}} \\
\hline
CNN-Small & 0.653 & 0.783 & 0.536 & 0.039 & 0.071 \\
CNN-Medium & 0.656 & 0.768 & 0.528 & 0.038 & 0.070 \\
CNN-Large & 0.642 & 0.749 & 0.528 & 0.038 & 0.070 \\
\hline
\multicolumn{6}{c}{\textbf{Transformer Family}} \\
\hline
Transformer-Medium & 0.632 & 0.778 & 0.447 & 0.033 & 0.060 \\
Transformer-Large & 0.631 & 0.764 & 0.485 & 0.036 & 0.066 \\
\hline
\end{tabular}
}
\end{table}
In contrast to the Sanity Proxy, the Musical Proxy results show minimal sensitivity to capacity scaling. Within the CNN family, NDCG@5 increases slightly from the small to the medium model and then decreases at the large tier, forming a non-monotonic pattern with small absolute differences. Within the Transformer family, the medium and large models exhibit nearly identical NDCG@5 values, indicating negligible impact of additional capacity. Across both families, performance differences remain small, suggesting that capacity scaling does not substantially affect ranking quality when relevance is defined by abstract semantic proxies.

\subsection{Capacity--Performance Trends Within Architectural Families}

This section analyzes scaling behavior strictly within architectural families to isolate the effect of parameter count from architectural inductive bias.

\subsubsection{CNN Scaling Behavior}

Within the CNN family, scaling behavior differs by task. On the Sanity Proxy, NDCG@5 shows a slight decrease from the small to the medium tier, followed by an increase at the large tier, indicating a non-monotonic pattern. On the Musical Proxy, scaling effects are weaker: NDCG@5 increases marginally from the small to the medium model and decreases at the large tier, with small absolute differences across models. These results indicate that scaling CNN-based models does not consistently improve retrieval performance and may even degrade it depending on the task.

\subsubsection{Transformer Scaling Behavior}

Within the Transformer family, scaling effects are also task-dependent. On the Sanity Proxy, performance decreases from the medium to the large model across all metrics, indicating that increased capacity does not lead to improved ranking quality in this setting. On the Musical Proxy, the medium and large models perform similarly, showing negligible differences in NDCG@5. These results indicate that scaling Transformer-based models does not consistently improve retrieval performance under either task.

\subsection{Cross-Family Descriptive Comparison}

Comparing performance across families reveals that scaling effects depend more on task structure than on model capacity alone. On the Sanity Proxy, both CNN and Transformer families exhibit non-monotonic behavior, with improvements observed only at specific capacity tiers. On the Musical Proxy, neither family shows consistent performance gains from increased capacity, and differences remain small across models. These comparisons are descriptive and do not imply causal superiority of one architecture over another, as architecture, parameter count, and pretraining regime co-vary.

\subsection{Embedding Extraction Cost}

Latency is measured as the average processing time per track in the embedding extraction pipeline. 
All latency benchmarks were conducted on an Apple Mac Mini M1 (2020) with 16GB RAM, using sequential inference on the CPU without GPU acceleration and a batch size of 1.

To complement the ranking results, we also report embedding extraction latency and model size for each model during offline feature generation.

\begin{table}[!htbp]
\centering
\caption{Embedding extraction cost and model size across model capacities. Parameter counts for CNN-based models are obtained from \cite{kong2020panns}, while Transformer parameter counts are based on publicly available pretrained checkpoints.}
\label{tab:latency_results}
\begin{tabular}{lcc}
\hline
Model & Params (M) & Latency (ms/track) \\
\hline
CNN-Small (Cnn6) & 4.8 & 2178.87 \\
CNN-Medium (Cnn10) & 5.2 & 3108.93 \\
CNN-Large (Cnn14) & 80.8 & 4217.94 \\
Transformer-Medium & 95 & 23145.73 \\
Transformer-Large & 330 & 55724.24 \\
\hline
\end{tabular}
\end{table}

Table~\ref{tab:latency_results} shows that embedding extraction latency increases substantially with model capacity. Within the CNN family, latency rises from 2{,}179\,ms per track (Cnn6, 4.8\,M parameters) to 4{,}218\,ms (Cnn14, 80.8\,M), a roughly 2$\times$ increase for a 17$\times$ growth in parameter count. The cost is more pronounced for Transformer models: latency increases from 23{,}146\,ms (MERT-95M) to 55{,}724\,ms (MERT-330M), a 2.4$\times$ increase for a 3.5$\times$ increase in parameters. These costs correspond to offline catalog ingestion, not online query time; once embeddings are indexed, retrieval is performed via similarity search and is independent of the embedding model. Nevertheless, the ingestion cost is relevant for archival systems that must re-index on catalog updates or operate without dedicated GPU infrastructure.


\subsection{Summary of Empirical Findings}

The experimental results yield three main observations. First, increasing model capacity does not produce consistent performance improvements. Across both architectural families, scaling exhibits non-monotonic behavior, and in some cases, performance decreases at higher capacity levels.  Second, scaling behavior is strongly task-dependent. Structured composer retrieval shows larger variation across capacity tiers, whereas abstract character retrieval exhibits minimal sensitivity to model size. Third, evaluation metrics behave differently when the number of relevant items per query is large relative to the evaluation cutoff. While NDCG@5 and Precision@5 remain relatively stable across models, Recall@5 and F1@5 are consistently low because only a small fraction of the relevant items can be retrieved within a limited cutoff (K=5). Collectively, these findings highlight the role of task structure and evaluation conditions in shaping observed performance patterns.



% Section 5: Discussion
\section{Discussion}

\subsection{Capacity Effects and Conditions for Improvement}

The results show that embedding-model capacity does not produce consistent improvements in candidate ranking quality in this setting. The effect is task-dependent: for structured retrieval based on composer identity, capacity differences produce measurable variation in ranking quality, although improvements are not monotonic; for abstract semantic retrieval based on musical character, scaling provides little additional discrimination and performance differences remain small across capacity tiers. This observation is consistent with prior work showing that strong MIR performance does not necessarily translate into improved recommendation quality \cite{tamm2024}. In addition, the stylistic coherence of classical music within a single archive may reduce embedding-space diversity, further limiting the benefit of additional model capacity.

\subsection{Implications for Recommender System Design}

The cost of deploying a higher-capacity model in this setting is concrete and immediate: as reported in Section~\ref{sec:results}, embedding extraction latency scales from approximately 2{,}200\,ms per track for the smallest CNN model to over 55{,}000\,ms for the largest Transformer, a roughly 25-fold increase. The corresponding gain in ranking quality, however, is neither consistent nor monotonic across model families and tasks. Taken together, these results indicate diminishing returns from capacity scaling: computational cost increases substantially while ranking quality does not, making mid-sized models a more practical choice for catalog ingestion and similarity-based retrieval in this setting.

For practitioners designing content-based candidate generation pipelines under similar constraints, these findings suggest that task formulation and relevance signal quality often outweigh the effect of model capacity. When relevance is broadly distributed across many items---as in the musical character proxy---scaling model size does not yield meaningful improvements. These results caution against treating model scaling as a default first step in system improvement under data-constrained conditions.

\subsection{Limitations}

This study has several limitations. First, the dataset is small and restricted to a single-domain classical music archive, limiting generalizability to larger and more diverse recommendation settings. However, this setting reflects realistic conditions in curated archival systems where content is limited and user interaction data is sparse. Although the dataset is small (N=203), it serves as a controlled and high-density setting for probing embedding space structure. In such a curated environment, where items are stylistically similar and semantically close, effective representations are expected to separate fine-grained distinctions. If capacity scaling does not improve separability under these conditions, it is unlikely to resolve semantic overlap in larger, more heterogeneous collections. Rather than modeling full-catalog retrieval at scale, this study isolates the item-to-item similarity ranking component. This component serves as a core building block in cold-start recommendation pipelines. Thus, the experimental setup should be interpreted as evaluating a local ranking stage rather than a full-catalog retrieval system. Second, the proxy relevance labels used in the evaluation are coarse and partially subjective, which constrains the assessment of fine-grained semantic similarity. Third, the evaluation is based on item-to-item retrieval rather than user-level recommendation, and therefore isolates candidate generation quality without capturing downstream ranking or personalization effects. Finally, the pseudo-labels used in the primary proxy task are partially derived from a pretrained emotion model, which may introduce bias. However, because these labels are fixed and applied consistently across all models, they primarily affect absolute performance values rather than the relative comparison of model capacity, which is the focus of this study. 

\subsection{Directions for Future Work}

Future work should validate these findings on larger and more diverse datasets. Incorporating richer semantic labels may improve the evaluation of abstract similarity in music retrieval. Finally, integrating user interaction signals and evaluating how embedding representations perform in user-level recommendation pipelines will be essential for understanding their practical impact in downstream settings.


% Section 6: Conclusion
\section{Conclusion}

This study investigated the effect of embedding-model capacity on candidate ranking quality in a content-based music retrieval setting for a small-scale classical music archive. Across both CNN and Transformer architectures, we find that increasing model capacity does not consistently improve ranking performance. Scaling behavior is non-monotonic, and in some cases, larger models yield degraded performance. Furthermore, the effect of capacity is strongly task-dependent. Structured retrieval tasks based on composer identity exhibit greater sensitivity to capacity differences, whereas abstract semantic retrieval tasks show minimal variation across model sizes. 

In addition, evaluation metrics behave differently when the number of relevant items is large relative to the evaluation cutoff. While ranking-based metrics such as NDCG@5 and Precision@5 remain relatively stable, set-based metrics such as Recall@5 and F1@5 are consistently low despite reasonable ranking quality. This highlights the importance of selecting evaluation metrics that align with the relevance distribution of the retrieval setting. Overall, this work suggests that, in small-scale archival recommendation settings, larger models incur substantially higher embedding extraction cost without delivering consistent ranking improvements, and that effective system design should favor lightweight or mid-sized models that provide comparable ranking quality at substantially lower computational cost, rather than defaulting to larger models.

% Manual bibliography workflow:
% Edit RecSys_refs.tex directly, then compile RecSys.tex twice with pdflatex.
\input{RecSys_refs.tex}

\end{document}
